\section{Actual persistent point directions preserve the source tangent cone}

\begin{lemma}[The unchanged finite tangent cone has a genuine coherent polynomial model]
\label{hir:actual-source-cone-polynomial-model}
\uses{hir:actual-finite-initial-vertex-whole-generation,hir:actual-translated-initial-ray-family-isomorphisms}
\lean{TranscendentalNakamaye.Pending.HironakaActualPersistentPointConeDerivatives.coneSheaf_isCoherent,TranscendentalNakamaye.Pending.HironakaActualPersistentPointConeDerivatives.coneTranslatedIdeal_eq,TranscendentalNakamaye.Pending.HironakaActualPersistentPointConeDerivatives.coneSheaf_translated_origin,TranscendentalNakamaye.Pending.HironakaActualPersistentPointConeDerivatives.cone_origin_formal}
\leanok
For every original analytic ideal, the actual finite homogeneous initial vertex
polynomials define a coherent holomorphic ideal on the entire coordinate space.
Its whole translated origin ideal is the literal translated initial polynomial
ideal at every point. Its formal origin extension is exactly the entire
homogeneous initial ideal of the original formal extension.
\end{lemma}
\begin{proof}
Use the genuine finite initial vertex polynomials of the original Taylor
extension, evaluated in the actual complex coordinates. Finite polynomial
functions are holomorphic everywhere, and their original finite germ span
is coherent with these same generators. Translating each polynomial at a
point gives exactly evaluation at the coordinate germ plus that point's
coordinate. Finite ideal spans therefore agree in both directions, rather
than only after taking their zero sets. Taylor expansion sends each
polynomial germ to its literal polynomial inclusion. The finite initial
vertex generation theorem identifies the resulting entire formal ideal
with the actual homogeneous initial ideal.
\end{proof}

\begin{lemma}[Actual point persistence makes the entire source cone constant along its direction]
\label{hir:actual-persistent-point-cone-ray}
\uses{hir:actual-source-cone-polynomial-model,hir:actual-analytic-point-normalized-bennett-all-orders,hir:actual-translated-initial-ray-family-isomorphisms}
\lean{TranscendentalNakamaye.Pending.HironakaActualPersistentPointConeDerivatives.translated_cone_hs_eq_of_zero_point_persistence}
\leanok
If the whole normalized transform in the zero-direction point chart retains
every Hilbert--Samuel value of the original analytic ideal, then the same
source cone has that entire profile at every point of the exceptional
coordinate line, including its origin.
\end{lemma}
\begin{proof}
The actual Bennett comparison factors through the dehomogenized source
initial ideal, which is genuinely isomorphic to the source cone translated
to the coordinate unit point. Persistence makes both ends of this chain of
inequalities equal, hence gives equality at that unit point at every order.
Every nonzero scalar on the line induces an actual germ algebra isomorphism
between the corresponding cone stalk and that unit stalk: each homogeneous
generator acquires only its own nonzero scalar power. At scalar zero, the
whole initial-ideal Hilbert--Samuel equality identifies the cone profile
with the original profile. These are all scalars, so the same entire
profile is constant on the actual whole line.
\end{proof}

\begin{theorem}[Actual zero-direction point persistence forces the literal source tangent derivative]
\label{hir:actual-persistent-point-source-cone-pderiv}
\uses{hir:actual-persistent-point-cone-ray,hir:actual-supplied-flattening-family-coupled-normal-data}
\lean{TranscendentalNakamaye.Pending.HironakaActualPersistentPointConeDerivatives.initial_pderiv_of_zero_point_hs_persistence}
\leanok
For any original analytic ideal and any coordinate pivot, actual all-order
Hilbert--Samuel persistence of its zero-direction point-normalized transform
forces that coordinate partial derivative to preserve the ENTIRE homogeneous
initial formal ideal of the SAME original Taylor extension. No cone
constancy, chart, rank, derivative-stability or presentation certificate is
an input.
\end{theorem}
\begin{proof}
Permute the pivot to the first coordinate and use the unchanged finite
homogeneous cone polynomials in those actual coordinates. The preceding
whole stalk equalities give constant actual Macaulay ranks along the first
coordinate line at every finite order. The genuine constant-rank family
producer constructs actual pivot neighbourhoods and first-symbol derivative
closure, and hence partial-derivative closure of the homogeneous initial
ideal of this polynomial-family origin ideal.

An ideal generated by homogeneous polynomials is its own homogeneous
initial ideal. Indeed, each homogeneous component of a finite sum of
coefficient-times-generator products remains in that ideal: a product's
component is the complementary component of its coefficient times the same
homogeneous generator, or is zero below that generator's degree. Every
initial form therefore lies in the ideal, while each nonzero homogeneous
generator is its own initial form. Thus the just obtained derivative acts
on the literal permuted source cone. Polynomial differentiation commutes
with the actual coordinate permutation. Apply this to every genuine finite
cone generator and undo the permutation. Finally linearity and the
Leibniz rule extend the conclusion to its entire formal ideal, including
arbitrary formal coefficient multiples. The empty-variable case has no
pivot and requires no dimension assumption. This gives a genuine input to
the algebraic own-order converse; it does not assert the remaining relative
or uniform joint-profile theorem.
\end{proof}
